%% kapa: exam=1 round=36 year=2025 period=2 form=A subject=appraisal_law subject_ko=감정평가관계법규 seq=35 number=35
%% source: Q-Net 감정평가사 기출문제 (cst00309) artlSeq=5231525 file=raw/1008/5231525/2025년 제36회 감정평가사 1차시험 2교시 문제지 원본.pdf
%% conversion: kordoc 4.24.51 -> tools/gen1.py (auto); review=proofread-sonnet; flags=none
\begin{question}{35}
부동산등기법령상 이의에 관한 설명으로 옳은 것은?
\begin{choices}
  \item 관할 지방법원은 이의신청에 대하여 결정한 후에 등기관에게 가등기 또는 이의가 있다는 뜻의 부기등기를 명령할 수 있다.
  \item 새로운 사실이나 새로운 증거방법을 근거로 이의신청을 할 수 있다.
  \item 등기관은 등기를 마치기 전에 이의가 이유 있다고 인정하면 3일 이내에 의견을 붙여 이의신청서 또는 이의신청정보를 지방법원에 보내야 한다.
  \item 등기관의 결정에 이의를 한 때에는 관할 지방법원의 결정때까지 집행이 정지된다.
  \item 관할 지방법원은 이의가 이유 있다고 인정하면 등기관에게 그에 해당하는 처분을 명령하고 그 뜻을 이의신청인은 물론 등기상 이해관계 있는 자에게도 알려야 한다.
\end{choices}
\end{question}
\answer{5}
